\section{预定义样式}\label{custom}
如何进行取决于你如何使用该宏包。一个好方法是尽可能使用预定义样式，将内容与形式分离。如果你计划创建一个包含大量图形的文档（就像本手册一样），这是正确的方法。我们将看到如何为文档定义全局样式，以及如何在局部使用样式。

文件 \tkzname{tkz-euclide.cfg} 包含了主要对象的预定义样式，其中最重要的是点、直线、线段、圆、圆弧和圆规记号。
 如果你总是使用相同的样式并创建大量图形，那么值得创建自己的样式。为此，你需要知道哪些特性可以修改，以及 \TIKZ\ 的一些概念。
 
 预定义样式是全局的：它们在图形创建之前就已存在，最好不要在两个图形之间修改它们。另一方面，它们可以在一个图形内部被临时修改；此时它们被局部定义，不会影响其他图形。
 
 下面是我如何定义你正在阅读的本文档的样式。

\begin{tkzltxexample}[]
  \tkzSetUpColors[background=white,text=black]  
  \tkzSetUpPoint[size=2,color=teal]
  \tkzSetUpLine[line width=.4pt,color=teal]
  \tkzSetUpCompass[color=orange, line width=.4pt,delta=10]
  \tkzSetUpArc[color=gray,line width=.4pt]
  \tkzSetUpStyle[orange]{new}
\end{tkzltxexample}

宏 \tkzcname{tkzSetUpColors} 设置背景色和文本颜色。如果你不使用它，则会使用文档的颜色和字体。下面我们来看看如何定义主要对象的样式。

\section{点样式}
点的定义方式如下：
\begin{tkzltxexample}[]
\tikzset{point style/.style = {%
       draw         = \tkz@euc@pointcolor,
       inner sep    = 0pt,
       shape        = \tkz@euc@pointshape,
       minimum size = \tkz@euc@pointsize,
       fill         = \tkz@euc@pointcolor}}  
\end{tkzltxexample}

你当然可以使用 \tkzcname{tikzset}，但宏包为此提供了一个宏：\tkzcname{tkzSetUpPoint}，它可以全局或局部使用。

\subsection{\tkzcname{tkzSetUpPoint} 的用法}

\begin{NewMacroBox}{tkzSetUpPoint}{\oarg{local options}}%
\begin{tabular}{lll}%
选项 &  默认值 & 说明                 \\ 
\midrule
\TOline{color}{black}{点的颜色} 
\TOline{size}{3}{点的大小} 
\TOline{fill}{black!50}{点的内部颜色} 
\TOline{shape}{circle}{点的形状：circle、cross 或 cross out} 
\end{tabular}
\end{NewMacroBox}



\subsubsection{全局样式或局部样式}
首先，下面是用本手册的样式创建的一个图形；然后在 \tkzNameEnv{tikzpicture} 环境内部修改了点的样式。

你可以全局或局部使用宏 \tkzcname{tkzSetUpPoint}。如果你将其放在导言区或第一个图形之前，点的样式将对文档中所有图形有效。你可以通过将该命令放在 \tkzNameEnv{tikzpicture} 环境内部来局部使用另一种样式。


\begin{tkzexample}[latex=7cm,small]
\begin{tikzpicture}
  \tkzDefPoints{0/0/A,5/0/B,3/2/C,3/1/D}
  \tkzDrawPolygon(A,B,C)
  \tkzDrawPoints(A,B,C)
  \tkzLabelPoints(A,B) 
  \tkzLabelPoints[above right](C) 
\end{tikzpicture}
\end{tkzexample}

\subsubsection{局部样式}
在第二个图形中局部修改了点的样式。
\begin{tkzexample}[latex=7cm,small]
\begin{tikzpicture}
  \tkzSetUpPoint[size=4,color=red,fill=red!20]
  \tkzDefPoints{0/0/A,5/0/B,3/2/C,3/1/D}
  \tkzDrawPolygon(A,B,C)
  \tkzDrawPoints(A,B,C)
  \tkzDrawPoint[shape=cross out,thick](D)
  \tkzLabelPoints(A,B) 
  \tkzLabelPoints[above right](C) 
\end{tikzpicture}
\end{tkzexample}

\subsubsection{\tkzname{Style} 与 \tkzname{scope}}
点恢复了初始样式。点 D 具有一个受环境 \tkzNameEnv{scope} 限制的新样式。也可以使用 |{...}| 或点恢复了初始样式。点 $D$ 具有一个受环境 \tkzNameEnv{scope} 限制的新样式。也可以使用 |{...}| 或 |\begingoup  ... \endgroup|。

\begin{tkzexample}[latex=7cm,small]
\begin{tikzpicture}
  \tkzDefPoints{0/0/A,5/0/B,3/2/C,3/1/D}
  \tkzDrawPolygon(A,B,C)
  \tkzDrawPoints(A,B,C)
  \begin{scope}
    \tkzSetUpPoint[size=4,color=red,fill=red!20]
    \tkzDrawPoint(D)
  \end{scope}
  \tkzLabelPoints(A,B) 
  \tkzLabelPoints[above right](C) 
\end{tikzpicture}
\end{tkzexample}

\subsubsection{\tkzcname{tkzSetUpPoint} 简单示例}

\begin{tkzexample}[latex=5cm,small]
\begin{tikzpicture}
  \tkzSetUpPoint[shape = cross out,color=blue]
  \tkzDefPoint(2,1){A}
  \tkzDefPoint(4,0){B}
  \tkzDrawLine(A,B)
  \tkzDrawPoints(A,B)
\end{tikzpicture}
\end{tkzexample}

\subsubsection{在 group 内部使用 \tkzcname{tkzSetUpPoint}}
\begin{tkzexample}[latex=8cm,small]
\begin{tikzpicture}
  \tkzDefPoints{0/0/A,2/4/B,4/0/C,3/2/D}
  \tkzDrawSegments(A,B A,C A,D)
  {\tkzSetUpPoint[shape=cross out,
            fill= teal!50,
            size=4,color=teal]
  \tkzDrawPoints(A,B)}
  \tkzSetUpPoint[fill= teal!50,size=4,
               color=teal]
   \tkzDrawPoints(C,D)
  \tkzLabelPoints(A,B,C,D)
\end{tikzpicture}
\end{tkzexample}

\section{直线样式}

有几种方法可以改变直线的样式：使用 \tkzcname{tkzSetUpLine}，或者直接使用\\ |\tikzset{line style/.style = ... }|。

关于 \tkzname{line width} 的提醒：有许多预定义样式可以更"自然"地设置线宽。你也可以重新定义这些样式。


\medskip
\begin{tabular}{cc}
预定义样式 & \tkzname{line width} 的值 \cr
\midrule
ultra thin    &  0.1 pt \cr
very thin     &  0.2 pt \cr
thin          &  0.4 pt \cr
semithick     &  0.6 pt \cr
thick         &  0.8 pt \cr
very thick    &  1.2 pt \cr
ultra thick   &  1.6 pt \cr
\midrule
\end{tabular}


\subsection{\tkzcname{tkzSetUpLine} 的用法} \label{tkzsetupline}
该宏定义所有直线的样式。

\begin{NewMacroBox}{tkzSetUpLine}{\oarg{local options}}%
\begin{tabular}{lll}%
选项 &  默认值 & 说明                 \\
\midrule
\TOline{color}{black}{构造线的颜色}
\TOline{line width}{0.4pt}{构造线的粗细}
\TOline{style}{solid}{构造线的样式}
\TOline{add}{.2 and .2}{改变线段的长度}
\end{tabular}
\end{NewMacroBox}

\subsubsection{改变线宽}
\begin{tkzexample}[latex=8cm,small]
\begin{tikzpicture}[scale=.75]
\tkzSetUpLine[line width=1pt]
\begin{scope}[rotate=-90]
    \tkzDefPoints{0/6/A,10/0/B,10/6/C}
    \tkzDefPointBy[projection = onto B--A](C)
    \tkzGetPoint{H}
    \tkzMarkRightAngle[size=.4,
                       fill=teal!20](B,C,A)
    \tkzMarkRightAngle[size=.4,
                       fill=orange!20](B,H,C)
    \tkzDrawPolygon(A,B,C)
    \tkzDrawSegment[new](C,H)
\end{scope}
 \tkzLabelSegment[below](C,B){$a$}
 \tkzLabelSegment[right](A,C){$b$}
 \tkzLabelSegment[left](A,B){$c$}
 \tkzLabelSegment[color=red](C,H){$h$}
 \tkzDrawPoints(A,B,C)
 \tkzLabelPoints[above left](H)
 \tkzLabelPoints(B,C)
 \tkzLabelPoints[above](A)
\end{tikzpicture}
\end{tkzexample}

\subsubsection{改变直线样式}

\begin{tkzexample}[latex=7cm,small]
\begin{tikzpicture}[scale=.5]
\tikzset{line style/.style = {color = gray,
                             style=dashed}}
\tkzDefPoints{1/0/A,4/0/B,1/1/C,5/1/D}
\tkzDefPoints{1/2/E,6/2/F,0/4/A',3/4/B'}
\tkzCalcLength(C,D)
\tkzGetLength{rCD}
\tkzCalcLength(E,F)
\tkzGetLength{rEF}
\tkzInterCC[R](A',\rCD)(B',\rEF)
\tkzGetPoints{I}{J}
\tkzDrawLine(A',B')
\tkzCompass(A',B')
\tkzDrawSegments(A,B C,D E,F)
\tkzDefCircle[R](A',\rCD) \tkzGetPoint{a'}
\tkzDefCircle[R](B',\rEF)\tkzGetPoint{b'}
\tkzDrawCircles(A',a' B',b')
\begin{scope}
  \tkzSetUpLine[color=red]
  \tkzDrawSegments(A',I B',I)
\end{scope}
\tkzDrawPoints(A,B,C,D,E,F,A',B',I,J)
\tkzLabelPoints(A,B,C,D,E,F,A',B',I,J)
\end{tikzpicture}
\end{tkzexample}

\subsubsection{延长直线}
\begin{tkzexample}[latex=7cm,small]
  \begin{tikzpicture}[scale=.75]
  \tkzSetUpLine[add=.5 and .5]
  \tkzDefPoints{0/0/A,4/0/B,1/3/C}
  \tkzDrawLines(A,B B,C A,C)
  \tkzDrawPolygon[red,thick](A,B,C)
  \tkzSetUpPoint[size=4,circle,color=red,fill=red!20]
  \tkzDrawPoints(A,B,C)
  \end{tikzpicture}
\end{tkzexample}

\section{圆弧样式}

\subsection{宏 \tkzcname{tkzSetUpArc}}
\begin{NewMacroBox}{tkzSetUpArc}{\oarg{local options}}%
\begin{tabular}{lll}%
选项 &  默认值 & 说明                 \\
\midrule
\TOline{color}{black}{直线的颜色}
\TOline{line width}{0.4pt}{直线的粗细}
\TOline{style}{solid}{构造线的样式}
\TOline{delta}{0}{在每端将圆弧延长 \tkzname{delta} 度}
\end{tabular}

\medskip
自第 5.14 版起，用本宏设置的 \tkzname{delta} 仅作用于 \tkzcname{tkzDrawArc}、\tkzcname{tkzMarkArc} 和 \tkzcname{tkzLabelArc} 所绘制的圆弧；圆规记号有自己的 \tkzname{delta}，由 \tkzcname{tkzSetUpCompass} 设置。
\end{NewMacroBox}

\subsubsection{\tkzcname{tkzSetUpArc} 的用法}
\begin{tkzexample}[latex=7cm,small]
\begin{tikzpicture}
\def\r{3} \def\angle{200}
\tkzSetUpArc[delta=10,color=purple,line width=.2pt]
\tkzSetUpLabel[font=\scriptsize,red]
\tkzDefPoint(0,0){O}
\tkzDefPoint(\angle:\r){A}
\tkzInterCC(O,A)(A,O) \tkzGetPoints{C'}{C}
\tkzInterCC(O,A)(C,O) \tkzGetPoints{D'}{D}
\tkzInterCC(O,A)(D,O) \tkzGetPoints{X'}{X}
\tkzDrawCircle(O,A)
\tkzDrawArc(A,C')(C)
\tkzDrawArc(C,O)(D)
\tkzDrawArc(D,O)(X)
\tkzDrawLine[add=.1 and .1](A,X)
\tkzDrawPoints(O,A)
\tkzSetUpPoint[size=3,color=purple,fill=purple!10]
\tkzDrawPoints(C,C',D,X)
\tkzLabelPoints[below left](O,A)
\tkzLabelPoints[below](C')
\tkzLabelPoints[below right](X)
\tkzLabelPoints[above](C,D)
\end{tikzpicture}
\end{tkzexample}

\section{圆规样式}
圆规记号有助于理解图形的构造，通过显示获得某些点所需的圆弧。

\subsection{宏 \tkzcname{tkzSetUpCompass} }
\begin{NewMacroBox}{tkzSetUpCompass}{\oarg{local options}}%
\begin{tabular}{lll}%
选项 &  默认值 & 说明                 \\
\midrule
\TOline{color}{black}{构造线的颜色}
\TOline{line width}{0.4pt}{构造线的粗细}
\TOline{style}{solid}{直线样式：solid、dashed、dotted、...}
\TOline{delta}{0}{改变圆规记号的长度}
\end{tabular}
\end{NewMacroBox}

\subsubsection{\tkzcname{tkzSetUpCompass} 的用法}

\begin{tkzexample}[latex=7cm,small]
\begin{tikzpicture}
  \tkzSetUpCompass[color=red,delta=15]
  \tkzDefPoint(1,1){A}
  \tkzDefPoint(6,1){B}
  \tkzInterCC[R](A,4)(B,4) \tkzGetPoints{C}{D}
  \tkzCompass(A,C)
  \tkzCompass(B,C)
  \tkzDrawPolygon(A,B,C)
  \tkzDrawPoints(A,B,C)
\end{tikzpicture} 
\end{tkzexample}

\subsubsection{\tkzcname{tkzSetUpCompass} 与 \tkzcname{tkzShowLine} 配合使用}

\begin{tkzexample}[latex=7cm,small]
\begin{tikzpicture}[scale=.75]
\tkzSetUpStyle[bisector,size=2,gap=3]{showbi}
\tkzSetUpCompass[color=teal,line width=.3 pt] 
\tkzDefPoints{0/1/A, 8/3/B, 3/6/C} 
\tkzDrawPolygon(A,B,C) 
\tkzDefLine[bisector](B,A,C) \tkzGetPoint{a} 
\tkzDefLine[bisector](C,B,A) \tkzGetPoint{b} 
\tkzShowLine[showbi](B,A,C) 
\tkzShowLine[showbi](C,B,A) 
\tkzInterLL(A,a)(B,b) \tkzGetPoint{I} 
\tkzDefPointBy[projection= onto A--B](I)
\tkzGetPoint{H} 
\tkzDrawCircle[new](I,H) 
\tkzDrawSegments[new](I,H) 
\tkzDrawLines[add=0 and .2,new](A,I B,I)
\end{tikzpicture}
\end{tkzexample} 





\section{标签样式}

\subsection{宏 \tkzcname{tkzSetUpLabel} }
宏 \tkzcname{tkzSetUpLabel} 定义点标签的样式。
\begin{NewMacroBox}{tkzSetUpStyle}{\oarg{local options}}%
  选项与 \TIKZ\ 的相同
\end{NewMacroBox}

\subsubsection{\tkzcname{tkzSetUpLabel} 的用法}
\begin{tkzexample}[latex=6cm,small]
\begin{tikzpicture}[scale=.75]
  \tkzSetUpLabel[font=\scriptsize,red]
  \tkzSetUpStyle[line width=1pt,teal]{XY}
  \tkzInit[xmin=-3,xmax=3,ymin=-3,ymax=3]
  \tkzDrawX[noticks,XY]
  \tkzDrawY[noticks,XY]
  \tkzDefPoints{1/0/A,0/1/B,-1/0/C,0/-1/D}
  \tkzDrawPoints[teal,fill=teal!30,size=6](A,...,D)
  \tkzLabelPoint[above right](A){$A(1,0)$}
  \tkzLabelPoint[above right](B){$B(0,1)$}
  \tkzLabelPoint[above left](C){$C(-1,0)$}
  \tkzLabelPoint[below left](D){$D(0,-1)$}
\end{tikzpicture}
\end{tkzexample}


\section{自定义样式}
你可以用 \tkzcname{tkzSetUpStyle} 定义自己的样式。

\subsection{宏 \tkzcname{tkzSetUpStyle} }
\begin{NewMacroBox}{tkzSetUpStyle}{\oarg{local options}}%
  选项与 \TIKZ\ 的相同
\end{NewMacroBox}

\subsubsection{\tkzcname{tkzSetUpStyle} 的用法}
\begin{tkzexample}[latex=2cm,small]
\begin{tikzpicture}
  \tkzSetUpStyle[color=blue!20!black,fill=blue!20]{mystyle}
  \tkzDefPoint(0,0){O}
  \tkzDefPoint(0,1){A}
  \tkzDrawPoints(O) % general style
  \tkzDrawPoints[mystyle,size=4](A) % my style
  \tkzLabelPoints(O,A)
\end{tikzpicture}
\end{tkzexample}

\section{如何使用 \tkzname{arrows}}

在某些国家，箭头用于表示直线平行，
或表示射线，例如角的边。

下面是一些如何放置这些箭头的示例。
\tkzname{tkz-euclide} 会加载库 \tkzname{arrows.meta}。
 
|\usetikzlibrary{arrows.meta}|

该库提供不同样式的箭头；下例使用其中的一些。

\subsection{在直线、射线或线段端点处的箭头}
\tkzname{Stealth}、\tkzname{Triangle}、\tkzname{To}、\tkzname{Latex}、\dots，可以与 \tkzname{reversed} 组合。在一个或两个端点处放置箭头非常容易。

\begin{enumerate}
\item \tkzname{-Triangle} 与 \tkzname{Segment}
 \begin{tkzexample}[latex=6cm,small]
    \begin{tikzpicture}
      \tkzDefPoints{0/0/A,4/0/B}
      \tkzDrawSegment[-Triangle](A,B)
    \end{tikzpicture}
  \end{tkzexample}
\item \tkzname{Stealth-Stealth} 与 \tkzname{Segment}
  \begin{tkzexample}[latex=6cm,small]
    \begin{tikzpicture}
      \tkzDefPoints{0/0/A,4/0/B}
      \tkzDrawSegment[Stealth-Stealth](A,B)
    \end{tikzpicture}
  \end{tkzexample}
\item \tkzname{Latex-Latex}   与 \tkzname{Line}
  \begin{tkzexample}[latex=6cm,small]
    \begin{tikzpicture}
      \tkzDefPoints{0/0/A,4/0/B}
      \tkzDrawLine[red,Latex-Latex](A,B)
      \tkzDrawPoints(A,B)
    \end{tikzpicture}
  \end{tkzexample}
\item \tkzname{To-To} 与 \tkzname{Segment}
  \begin{tkzexample}[latex=6cm,small]
    \begin{tikzpicture}
      \tkzDefPoints{0/0/A,4/0/B}
      \tkzDrawSegment[To-To](A,B)
    \end{tikzpicture}
  \end{tkzexample}
\item \tkzname{Latex-Latex} 与 \tkzname{Segment}
  \begin{tkzexample}[latex=6cm,small]
    \begin{tikzpicture}
      \tkzDefPoints{0/0/A,4/0/B}
      \tkzDrawSegment[Latex-Latex](A,B)
    \end{tikzpicture}
  \end{tkzexample}
\item \tkzname{Latex-}  与 \tkzname{Segment}
  \begin{tkzexample}[latex=6cm,small]
    \begin{tikzpicture}
      \tkzDefPoints{0/0/A,4/0/B}
      \tkzDrawSegment[Latex-](A,B)
    \end{tikzpicture}
  \end{tkzexample}
\item \tkzname{-Latex}  与 \tkzname{Segments}
  \begin{tkzexample}[latex=6cm,small]
\begin{tikzpicture}
 \tkzDefPoints{0/0/A,4/0/B,5/-2/C}
 \tkzDrawSegments[-Latex](A,B A,C)
\end{tikzpicture}
\end{tkzexample}
\end{enumerate}

\subsubsection{缩放箭头头部}

\begin{tkzexample}[latex=6cm,small]
\begin{tikzpicture}
 \tkzDefPoints{0/0/A,4/0/B}
 \tkzDrawSegment[{Latex[scale=2]}-{Latex[scale=2]}](A,B)
\end{tikzpicture}
\end{tkzexample}

\subsubsection{使用 vector 样式}
|\tikzset{vector style/.style={>=Latex,->}}|

你可以重新定义这个样式。
\begin{tkzexample}[latex=6cm,small]
\begin{tikzpicture}
 \tkzDefPoints{0/0/A,4/0/B}
\tkzDrawSegment[vector style](A,B)
\end{tikzpicture}
\end{tkzexample}
  
\subsection{在线段中点处的箭头}

直线上的箭头表示这些直线平行。这取决于国家：在法国，我们更倾向于在图形外部写出 $(AB) \parallel (DC)$。以下代码改编自 \tkzname{Muzimuzhi Z} 在 \href{https://tex.stackexchange.com/questions/632596/how-to-manage-argument-pattern-keys-and-subways}{tex.stackexchange.com} 上的回答。

\medskip
 语法：\\

 \begin{itemize}
\item |tkz arrow|（默认为 \tkzname{Latex}）
\item |tkz arrow=<arrow end tip>|
\item |tkz arrow=<arrow end tip> at <pos> (<pos> = .5 by default)|
\item |tkz arrow={<arrow end tip>[<arrow options>] at <pos>}| 可使用选项 \tkzname{scale}
 \end{itemize}

用法示例：\\

|\tkzDrawSegment[tkz arrow=Stealth] (A,B)|\\
|\tkzDrawSegment[tkz arrow={To[scale=3] at .4}](A,B)|\\
|\tkzDrawSegment[tkz arrow={Latex[scale=5,blue] at .6}](A,B)|

\subsubsection{在平行四边形中}
\begin{tkzexample}[latex=7cm,small]
\begin{tikzpicture}
 \tkzDefPoints{0/0/A,3/0/B,4/2/C} 
 \tkzDefParallelogram(A,B,C) 
 \tkzGetPoint{D}
 \tkzDrawSegments[tkz arrow](A,B D,C)
 \tkzDrawSegments(B,C D,A)
 \tkzLabelPoints(A,B) 
 \tkzLabelPoints[above right](C,D)
 \tkzDrawPoints(A,...,D)
\end{tikzpicture}
\end{tkzexample}

\subsubsection{与另一条直线平行的直线}
\begin{tkzexample}[latex=7cm,small]
\begin{tikzpicture}
 \tkzDefPoints{0/0/A,3/0/B,1/2/C} 
 \tkzDefPointWith[colinear= at C](A,B) 
 \tkzGetPoint{D}
 \tkzDrawSegments[tkz arrow=Triangle](A,B C,D)
 \tkzLabelPoints(A,B,C) 
 \tkzDrawPoints(A,...,C)
\end{tikzpicture}
\end{tkzexample}

\subsubsection{圆上的箭头}
你可以在圆的第一象限处放置一个箭头；旋转可以移动该箭头。
\begin{tkzexample}[latex=7cm,small]
\begin{tikzpicture}
\tkzDefPoints{0/0/A,3/0/B} 
\begin{scope}[rotate=150]
 \tkzDrawCircle[tkz arrow={Latex[scale=2,red]}](A,B)
\end{scope}
\end{tikzpicture}
\end{tkzexample}

\subsection{在多边形的所有边上放置箭头}
一些用户向我询问一种在多边形每条边上放置箭头的方法。我使用了 Paul Gaborit 在 \href{https://tex.stackexchange.com/questions/3161/tikz-how-to-draw-an-arrow-in-the-middle-of-the-line}{tex.stackexchange.com} 上提出的样式。

|\tikzset{tkz arrows/.style=|\\
|{postaction={on each path={tkz arrow={Latex[scale=2,color=black]}}}}}|   

你可以改变这个样式。使用 \tkzname{tkz arrows} 你可以在多边形的每条线段上放置一个箭头。

\subsubsection{用 \tkzname{tkz arrows} 在每条线段上放置箭头} 
\begin{tkzexample}[latex=7cm,small]
\begin{tikzpicture}
 \tkzDefPoints{0/0/A,3/0/B}  
 \tkzDefSquare(A,B) \tkzGetPoints{C}{D}
 \tkzDrawPolygon[tkz arrows](A,...,D)
\end{tikzpicture}
\end{tkzexample}

\subsubsection{将 \tkzname{tkz arrows} 用于圆}
\begin{tkzexample}[latex=7cm,small]
\begin{tikzpicture}
 \tkzDefPoints{0/0/A,3/0/B} 
 \tkzDrawCircle[tkz arrows](A,B)
\end{tikzpicture}
\end{tkzexample}

\endinput
